\subsection{Contador Síncrono MOD $6$}

Seja $D$ o estado atual do contador, restrito ao conjunto de valores válidos.

\paragraph{Incremento}
A operação de incremento é definida por:
\[
D \mapsto (D + 1) \bmod 6
\]

A Tabela~\ref{tab:inc_mod6} apresenta a função de transição.

\begin{table} [H] \label{tab:dec_mod6}
\begin{center}
\begin{tabular}{c|ccc}
    $D$ & $O$ & Carry \\
    \hline
    $000$ & $001$ & $0$ \\
    $001$ & $010$ & $0$ \\
    $010$ & $011$ & $0$ \\
    $011$ & $100$ & $0$ \\
    $100$ & $101$ & $0$ \\
    $101$ & $000$ & $1$ \\
    $110$ & $000$ & $0$ \\
    $111$ & $000$ & $0$ \\    
\end{tabular}
\end{center}
\end{table}
    

As expressões booleanas minimizadas são:

\begin{align*}
    O_2 &= (\lnot D_2 \land D_1 \land D_0) \lor (D_2 \land \lnot D_1 \land \lnot D_0) \\[4pt]
    O_1 &= (\lnot D_2 \land \lnot D_1 \land D_0) \lor (\lnot D_2 \land D_1 \land \lnot D_0) \\[4pt]
    O_0 &= (\lnot D_2 \land \lnot D_0) \lor (\lnot D_1 \land \lnot D_0) \\[4pt]
    Carry &= D_2 \land \lnot D_1 \land D_0
\end{align*}


O sinal \textit{Carry} indica overflow.

As Figuras~\ref{fig:incmod6.png} e~\ref{fig:incmod6-red.png} apresentam a implementação do circuito.

\imgbig{incmod6.png}{Circuito lógico para incremento em módulo 6.}
\imgbig{incmod6-red.png}{Implementação em redstone do circuito de incremento em módulo 6.}

\paragraph{Decremento}
A operação de decremento é definida por:
\[
D \mapsto (D - 1) \bmod 6
\]

A Tabela~\ref{tab:dec_mod6} apresenta a função de transição.

\begin{table} [H] \label{tab:dec_mod6}
\begin{center}
\begin{tabular}{c|ccc}    
    $D$ & $O$ & Carry \\
    \hline
    $000$ & $101$ & $0$ \\
    $001$ & $000$ & $0$ \\
    $010$ & $001$ & $0$ \\
    $011$ & $010$ & $0$ \\
    $100$ & $011$ & $0$ \\
    $101$ & $000$ & $1$ \\
    $110$ & $000$ & $0$ \\
    $111$ & $000$ & $0$ \\
\end{tabular}
\end{center}
\end{table}

\begin{align*}
    O_2 &= (\lnot D_2 \land \lnot D_1 \land \lnot D_0) \lor (D_2 \land D_0) \\[4pt]
    O_1 &= (D_1 \land D_0) \lor (D_2 \land \lnot D_1 \land \lnot D_0) \\[4pt]
    O_0 &= \lnot D_0 \\[4pt]
    Carry &= \lnot D_2 \land \lnot D_1 \land \lnot D_0
\end{align*}

O sinal \textit{Carry} indica underflow.

As Figuras~\ref{fig:decmod6.png} e~\ref{fig:decmod6-red.png} mostram a implementação do circuito.

\imgbig{decmod6.png}{Circuito lógico para decremento em módulo 6.}
\imgbig{decmod6-red.png}{Implementação em redstone do circuito de decremento em módulo 6.}